 %%%%%%%%%%%%%%%%%%%%%%%%%%%%%%%%%%%%%%%%%%%%%%%%%%%%%%%%%%%%%%%%%%%%%%%%%%%%%%%%
%2345678901234567890123456789012345678901234567890123456789012345678901234567890
%        1         2         3         4         5         6         7         8

\documentclass[a4paper, 10pt, conference]{ieeeconf}

\IEEEoverridecommandlockouts                              % This command is only needed if
 % you want to use the \thanks command

\overrideIEEEmargins                                      % Needed to meet printer requirements.

%In case you encounter the following error:
%Error 1010 The PDF file may be corrupt (unable to open PDF file) OR
%Error 1000 An error occurred while parsing a contents stream. Unable to analyze the PDF file.
%This is a known problem with pdfLaTeX conversion filter. The file cannot be opened with acrobat reader
%Please use one of the alternatives below to circumvent this error by uncommenting one or the other
%\pdfobjcompresslevel=0
%\pdfminorversion=4


% The following packages can be found on http:\\www.ctan.org
\usepackage[utf8]{inputenc}
\usepackage[T1]{fontenc}
\usepackage[ngerman]{babel}
\usepackage{graphicx} % for pdf, bitmapped graphics files
%\usepackage{epsfig} % for postscript graphics files
%\usepackage{mathptmx} % assumes new font selection scheme installed
%\usepackage{times} % assumes new font selection scheme installed
\usepackage{amsmath} % assumes amsmath package installed
%\usepackage{amssymb}  % assumes amsmath package installed

%%%%%%%%%%%%% added for project course control engineering
\usepackage[left=1.57cm,right=1.57cm,top=2.54cm,bottom=2.54cm]{geometry}
\usepackage{subfigure}
%%%%%%%%%%%%% added for project course control engineering


\usepackage{amssymb}
\usepackage{enumitem}
\usepackage{multirow}
\usepackage{amsthm}

%%%%%%%%%%%%% added for the Thermal Digital Twin (DT4TM) report %%%%%%%%%%%%%
\usepackage{booktabs}
\usepackage{tikz}
\usetikzlibrary{arrows.meta, positioning}
\usepackage{url} % \url{} breaks long filenames at _ . - / so they wrap inside narrow table columns (plain \texttt{} cannot break there)
\usepackage{tabularx} % X columns share \textwidth exactly and auto-wrap text,
                       % unlike plain tabular's c/l columns which size to content
                       % and silently overflow the page margin if content is wide
\usepackage[hidelinks]{hyperref} % makes \cite/\ref/\eqref clickable (jump to the
                       % bibliography/label); loaded last so it patches every
                       % preceding package's cross-referencing commands. `hidelinks`
                       % keeps links invisible (no colored text/boxes) so the printed
                       % look is unchanged -- only PDF click-navigation is added.
% ieeeconf.cls defines a FAKE \NAT@parse purely to make hyperref think natbib
% is loaded (so hyperref's natbib-compatibility code, not its plain-\cite code,
% handles \cite) -- but the class's own comment says this fools hyperref into
% NOT hyperlinking citation numbers at all. The class explicitly names this
% override as the fix (ieeeconf.cls, \NAT@parse definition): undefine it so
% \cite{...} actually links down to the matching \bibitem.
%%%%%%%%%%%%% added for the Thermal Digital Twin (DT4TM) report %%%%%%%%%%%%%

\theoremstyle{remark}
% Template default is \newtheorem{remark}{Remark}; renamed to the German
% "Anmerkung" only because this report is written in German. Same environment,
% same style -- no format change.
\newtheorem{remark}{Anmerkung}

% ieeeconf.cls hard-codes the English word "Appendix" for auto-generated
% appendix headings (see \@IEEEprocessthesectionargument); since this report
% is in German, override it here to "Anhang" so headings match the prose.
\makeatletter
\def\@IEEEprocessthesectionargument#1{%
\@ifmtarg{#1}{%
\@IEEEappendixsavesection*{Anhang \thesectiondis}%
\addcontentsline{toc}{section}{Anhang \thesection}}{%
\@IEEEappendixsavesection*{Anhang \thesectiondis \\* #1}%
\addcontentsline{toc}{section}{Anhang \thesection: #1}}}
\makeatother

\title{\LARGE \bf
Thermal Digital Twin für ein elektrodynamisches Levitationssystem
}
\author{ Dang Minh Hoang, Marina Borchert, Mohamed Aziz El Majid % <-this % stops a space
\thanks{Project Course Digital Twin, Summer Term 2026. Supervisors: Hr. Prof. Dr. rer. nat. Dirk Hartmann, Fr. Dr.-Ing. Melina Merkel}% project course: only change the supervisor(s)
}


\begin{document}



\maketitle
\thispagestyle{plain}
\pagestyle{plain}


%%%%%%%%%%%%%%%%%%%%%%%%%%%%%%%%%%%%%%%%%%%%%%%%%%%%%%%%%%%%%%%%%%%%%%%%%%%%%%%%
\begin{abstract}
Dieser Bericht beschreibt Konzept, Methodik und Softwarearchitektur eines
\emph{digitalen Zwillings} (Digital Twin) für die thermische Vorhersage am
elektrodynamischen Levitator nach dem TEAM-28-Benchmark (COMPUMAG). Statt
kommerzieller CAE-Software wurde ein selbst entwickelter Python-Löser
gewählt, weil das eigentliche Ziel keine einmalige Simulation, sondern ein
\emph{echtzeitfähiges} Vorhersagemodell ist, das sich schnell auf
intelligenten Endgeräten (Laptop, Smartphone) anzeigen lässt. Der Bericht
begründet diese Wahl physikalisch (die Joulschen Verluste skalieren
quadratisch mit dem Strom bei fester räumlicher Verteilung, was eine
einmalige Offline-FEM-Lösung plus echtzeitfähiges Reduced-Order-Model
erlaubt). Als Eingangsgröße dient dabei der \emph{gemessene} Spulenstrom
$I(t)$: Ein Hall-Stromsensor am Prüfstand speist den Zeitintegrator
fortlaufend, während die Temperatur die daraus berechnete Ausgangsgröße
ist. Der Bericht fasst bisherige sowie erwartete Ergebnisse zusammen und
unterscheidet dabei ausdrücklich zwischen dem bereits validierten
stationären Zustand und dem noch nicht validierten Zeitverlauf. Für die
Implementierung des physikalischen Modells wurde ein in die
Entwicklungsumgebung integriertes KI-Werkzeug genutzt, wobei alle
physikalischen Annahmen und Validierungskriterien vom Team definiert und
geprüft wurden.
\end{abstract}

%%%%%%%%%%%%%%%%%%%%%%%%%%%%%%%%%%%%%%%%%%%%%%%%%%%%%%%%%%%%%%%%%%%%%%%%%%%%%%%


\section{Einleitung}
\label{sec:intro}

\subsection{Die Aufgabenstellung}
Das TEAM-Problem~28~\cite{team28} ist ein etablierter Benchmark der
Computational-Electromagnetics-Gemeinschaft (COMPUMAG) zur Validierung
selbst geschriebener elektromagnetischer Löser: Eine leitfähige
Aluminiumscheibe levitiert über einer Spulenanordnung, die mit
Wechselstrom ($I_{\mathrm{rms}}\approx5\,\mathrm{A}$ gemessen,
$f=50\,\mathrm{Hz}$) betrieben wird. Der zeitlich veränderliche
magnetische Fluss induziert
Wirbelströme in der Scheibe, die sowohl eine Hubkraft (Levitation) als
auch eine Verlustwärme erzeugen. Der am Fachgebiet aufgebaute
Versuchsstand entspricht diesem Benchmark strukturell, besitzt
zusätzlich eine zweigeteilte Spulenanordnung (Innenspule ca.\
$1000$ Windungen, Außenspule ca.\ $500$ Windungen) mit
ferromagnetischem Kern ($\mu_r\approx1000$) und ferromagnetischem
Außenring, die selbst spürbar Wärme erzeugen.

\begin{figure}[t]
\centering
\begin{tikzpicture}[x=0.065cm, y=0.065cm,
    ironfill/.style={fill=gray!45},
    cufill/.style={fill=orange!35},
    alfill/.style={fill=cyan!25},
    lbl/.style={font=\scriptsize, align=center}]
  % Symmetrieachse
  \draw[dash dot] (0,15) -- (0,-58);
  \node[font=\tiny, right] at (2,16) {$r=0$};
  % r-Achse
  \draw[-{Stealth[length=4pt]}] (0,-58) -- (109,-58) node[below, font=\tiny] {$r$ [mm]};
  \draw (25.9,-58) -- (25.9,-56.5); \node[font=\tiny] at (25.9,-63) {25{,}9};
  \draw (61.9,-58) -- (61.9,-56.5); \node[font=\tiny] at (61.9,-63) {61{,}9};
  \draw (79.9,-58) -- (79.9,-56.5); \node[font=\tiny] at (79.9,-63) {79{,}9};
  \draw (102.9,-58) -- (102.9,-56.5); \node[font=\tiny] at (102.9,-63) {102{,}9};
  % Eisenkern
  \draw[ironfill] (0,0) rectangle (25.9,-53);
  \node[lbl] at (13,-26.5) {Eisenkern\\$\mu_r{=}1000$};
  % Innenspule
  \draw[cufill] (27.9,0) rectangle (61.9,-53);
  \node[lbl] at (44.9,-26.5) {Innenspule\\$N{=}1000$};
  % Eisenring
  \draw[ironfill] (64.9,0) rectangle (79.9,-53);
  \node[lbl, rotate=90] at (72.4,-26.5) {Eisenring};
  % Außenspule
  \draw[cufill] (82.9,0) rectangle (102.9,-53);
  \node[lbl] at (92.9,-26.5) {Außen-\\spule\\$N{=}500$};
  % Al-Scheibe (schematische Position -- Luftspalt umstritten, s. Abschnitt VIII)
  \draw[alfill] (0,6) rectangle (80,9);
  \node[lbl] at (40,7.5) {Al-Scheibe, $R{=}80$, $z{=}3$\,[mm]};
  \draw[<->, dashed] (85,0) -- (85,6);
  \node[font=\tiny] at (91,3) {$z_{\mathrm{gap}}$};
\end{tikzpicture}
\caption{Querschnitt der Gerätegeometrie, identisch zur Eingabe in
\texttt{params.yaml}. Eingangsgröße $I_{\mathrm{rms}}=5\,\mathrm A$:
rechnerisch $z_{\mathrm{gap}}\approx12\,\mathrm{mm}$, real beobachtet
$7$--$8\,\mathrm{mm}$ (Abschnitt~\ref{sec:limitations}).}
\label{fig:geometry}
\end{figure}

\begin{table}[t]
\centering\footnotesize
\caption{Reale,vermessene Gerätegeometrie für FEM-Löser (\texttt{params.yaml}).}
\label{tab:geometry}
\begin{tabular}{@{}lrrrr@{}}
\toprule
Bauteil & $r$\,[mm] & $\mu_r$ & $\sigma$\,[MS/m] & $N$ \\
\midrule
Eisenkern & 0{,}0--25{,}9 & 1000 & 1{,}0 & -- \\
Luftspalt & 25{,}9--27{,}9 & 1 & 0 & -- \\
Innenspule (Cu) & 27{,}9--61{,}9 & 1 & 59{,}6 & 1000 \\
Luftspalt & 61{,}9--64{,}9 & 1 & 0 & -- \\
Eisenring & 64{,}9--79{,}9 & 1000 & 1{,}0 & -- \\
Luftspalt & 79{,}9--82{,}9 & 1 & 0 & -- \\
Außenspule (Cu) & 82{,}9--102{,}9 & 1 & 59{,}6 & 500 \\
Al-Scheibe & 0{,}0--80{,}0 & 1 & 34{,}0 & -- \\
Luft $\to$ Rahmen & 102{,}9--130{,}4 & 1 & 0 & -- \\
Rahmenwand (Sperrholz) & 130{,}4--152{,}9 & 1 & 0 & -- \\
\bottomrule
\end{tabular}

\vspace{2pt}
{\footnotesize Höhe aller Kern-/Spulen-/Ringbereiche: $53\,\mathrm{mm}$;
Scheibendicke: $3\,\mathrm{mm}$ (oberhalb dieser Höhe positioniert, vgl.
Abbildung~\ref{fig:geometry}). Elektrische Anregung:
$I_{\mathrm{rms}}=5{,}0\,\mathrm A$ (Hauptbetriebspunkt, Variac-Dial
$220^\circ\approx195{,}5\,\mathrm V$) bzw.\ $7{,}8\,\mathrm A$
(Variac-Dial auf Maximalstellung $270^\circ$, $\approx240\,\mathrm V$,
Kalibrieranker), $f=50\,\mathrm{Hz}$.}
\end{table}

Abbildung~\ref{fig:geometry} und Tabelle~\ref{tab:geometry} dokumentieren
diese reale Geometrie vollständig und identisch zu \texttt{params.yaml} --
der \emph{einzigen} Quelle, aus der sowohl der eigene FEM-Löser als auch
ein potenzieller Vergleichslauf in kommerzieller Software gespeist werden.
Aus Transparenzgründen sei betont: Nur die Radien, die Windungszahlen und
die (per Magnettest bestätigte) Werkstoffklasse Eisen/Kupfer/Aluminium
sind direkt gemessen; Permeabilität $\mu_r$ und Leitfähigkeit $\sigma$
sind Referenz-/Literaturwerte für die jeweilige Werkstoffklasse
(Abschnitt~\ref{sec:limitations}), keine Messung an der konkreten Probe.

\subsection{Vom Simulationsmodell zum digitalen Zwilling}
Die Aufgabe besteht nicht darin, \emph{einmal} ein Temperaturfeld zu
berechnen, sondern ein Modell zu bauen, das (i) das Temperaturfeld
$T(r,z,t)$ \textbf{kontinuierlich} vorhersagt, angetrieben vom real
gemessenen Spulenstrom $I(t)$, (ii) gegen den realen Versuchsstand
\textbf{validiert} wird (IR-Kamera), und (iii) am Ende
über einen \textbf{QR-Code} auf einem Smartphone in Augmented Reality
betrachtet werden kann. Diese drei Anforderungen zusammen definieren
einen \emph{digitalen Zwilling} im eigentlichen Sinn~\cite{tao2019digitaltwin}
und sind der Ausgangspunkt aller im Folgenden begründeten methodischen
Entscheidungen.


\section{Physikalisches Grundkonzept}
\label{sec:concept}

Ein digitaler Zwilling, der bei jeder Stromänderung eine volle
FEM-Simulation neu rechnet, wäre auf einem Smartphone nicht realisierbar
(Rechenzeit: Sekunden bis Minuten pro Lauf). Die Lösung liegt in einer
physikalischen Beobachtung, nicht in einer Software-Abkürzung: Bei fester
Frequenz ($50\,\mathrm{Hz}$) und linearen Materialeigenschaften skalieren
alle Joulschen Verluste quadratisch mit dem Strom ($P\propto I^2$),
während sich die \textbf{räumliche Verteilung} der Verluste \emph{nicht}
ändert -- nur ihre Amplitude. Daraus folgt eine zweistufige Architektur:

\begin{enumerate}[nosep]
  \item \textbf{Offline (einmalig):} Die elektromagnetische FEM wird
        einmal bei einem Referenzstrom $I_{\mathrm{ref}}$
        gelöst und liefert eine ortsaufgelöste Verlustkarte
        $\hat q_{\mathrm{Al}}(r,z)$ sowie Spulenverluste
        $\hat P_{\mathrm{Cu}}$.
  \item \textbf{Online (Millisekunden):} Für einen beliebigen Strom $I$
        und die zuletzt bekannte Körpertemperatur $\bar T$ werden diese
        nur noch skalar multipliziert.
\end{enumerate}

Wichtig ist, dass die $\sigma(T)$-Korrektur für die beiden Körper
\textbf{gegenläufig} ist (Abschnitt~\ref{sec:physics}): Der
Wirbelstromverlust der Scheibe ist $\propto\sigma_{\mathrm{Al}}$, der
ohmsche Spulenverlust dagegen $\propto1/\sigma_{\mathrm{Cu}}$. Die
Echtzeitgleichungen lauten daher
\begin{align}
  q_{\mathrm{Al}}(r,z;I,\bar T) &= \hat q_{\mathrm{Al}}(r,z)
     \left(\frac{I}{I_{\mathrm{ref}}}\right)^{\!2}
     \frac{\sigma_{\mathrm{Al}}(\bar T)}{\sigma_{\mathrm{Al}}(T_{\mathrm{ref}})},
     \label{eq:realtime}\\[2pt]
  P_{\mathrm{Cu}}(I,\bar T) &= \hat P_{\mathrm{Cu}}
     \left(\frac{I}{I_{\mathrm{ref}}}\right)^{\!2}
     \frac{\sigma_{\mathrm{Cu}}(T_{\mathrm{ref}})}{\sigma_{\mathrm{Cu}}(\bar T)}.
     \label{eq:realtime_coil}
\end{align}
Beide speisen ein reduziertes thermisches Modell (ROM), das ebenfalls nur
mit Vektor-/Skalar-Arithmetik arbeitet.

\begin{remark}[Implementierungsstand]
\label{rem:sigma}
Gl.~\eqref{eq:realtime} ist im Laufzeitmodell vollständig umgesetzt.
Gl.~\eqref{eq:realtime_coil} ist es \emph{nicht}: Die Spulen- und
Eisenknoten werden zur Laufzeit nur mit $(I/I_{\mathrm{ref}})^2$
skaliert. Die Temperaturabhängigkeit der Spule steckt stattdessen
implizit in den gegen Messdaten gefitteten Konvektionsleitwerten $hA$
(Abschnitt~\ref{sec:validation}). Am Kalibrierpunkt ist das exakt, für
weit davon entfernte Betriebspunkte eine bekannte Vereinfachung.
\end{remark}

Diese Zerlegung ist der methodische Kern des Projekts: Die
"`Intelligenz"' steckt vollständig im einmaligen Offline-Löser; die
Echtzeitfähigkeit entsteht aus der physikalischen Struktur des Problems
(Linearität, feste Frequenz), nicht aus vereinfachter Physik.


\section{Physikalisches Modell}
\label{sec:physics}

Das Modell besteht aus drei gekoppelten, achsensymmetrischen
Teilproblemen, die in 2D $(r,z)$ gelöst und für die Darstellung zu 3D
rotiert werden. Die vollständige Herleitung befindet sich in
Anhang~\ref{app:derivation}.

\textbf{Elektromagnetik (Phasor):} Da nur die zyklusgemittelte
Verlustwärme für die (im Vergleich zu $50\,\mathrm{Hz}$ langsame) Thermik
relevant ist, wird das Feld nicht zeitschrittweise simuliert, sondern als
komplexer Phasor $A_\varphi(r,z)$ gelöst -- die hier verwendete
Formulierung über das magnetische Vektorpotential folgt der klassischen
Arbeit von Bíró und Preis~\cite{biro1989vectorpotential},
\begin{equation}
  -\nabla\!\cdot\!\left(\nu\nabla A_\varphi\right)
  + \nu\frac{A_\varphi}{r^{2}} + j\omega\sigma A_\varphi = J_s,
  \label{eq:em}
\end{equation}
mit $\nu=1/(\mu_r\mu_0)$, $\omega=2\pi f$ und der eingeprägten
Stromdichte $J_s$ in den Spulen (Innen-/Außenspule mit entgegengesetztem
Wicklungssinn). Im ferromagnetischen Kern/Ring fällt $\nu$ stark ab und
bündelt so den Fluss.

\textbf{Joulsche Verlustwärme:} Aus der induzierten Wirbelstromdichte
$J_e=-j\omega\sigma A_\varphi$ folgt die zyklusgemittelte Verlustdichte
\begin{equation}
  q(r,z) = \frac{|J_e|^{2}}{2\sigma}
         = \tfrac12\,\sigma\,\omega^{2}\,|A_\varphi|^{2}\ \;[\mathrm{W/m^3}],
  \label{eq:loss}
\end{equation}
deren quadratische Abhängigkeit von $|A_\varphi|$ die Grundlage der
$I^2$-Skalierung aus Gl.~\eqref{eq:realtime} ist.

\textbf{Temperaturabhängige Leitfähigkeit:} Da $70\,\mathrm K$ Erwärmung
den spezifischen Widerstand um $\approx27\,\%$ ändern, gilt
$\sigma(T)=\sigma_0/(1+\alpha(T-T_0))$ mit
$\alpha\approx3{,}9\times10^{-3}\,\mathrm{K^{-1}}$ -- mit gegenläufiger
Wirkung: Scheibenverluste $\propto\sigma_{\mathrm{Al}}$ (negative
Rückkopplung), Spulenverluste $\propto1/\sigma_{\mathrm{Cu}}$ (positive
Rückkopplung) -- siehe Gl.~\eqref{eq:realtime} und
\eqref{eq:realtime_coil} sowie Anmerkung~\ref{rem:sigma} zum
Umsetzungsstand. Beide
Effekte ändern nur die Amplitude, sodass die Zweistufenarchitektur gültig
bleibt.

\textbf{Wärmeleitung:} Die Verlustdichte $q$ speist die instationäre
Fourier-Gleichung
\begin{equation}
  \rho c_p\,\frac{\partial T}{\partial t} = \nabla\!\cdot\!(k\nabla T) + q,
  \label{eq:fourier}
\end{equation}
mit Robin-Randbedingung $-k\,\partial T/\partial n=h(T-T_\infty)$. Im
stationären Fall muss die Energiebilanz
$\int_V q\,\mathrm dV=\oint_\Gamma h(T-T_\infty)\,\mathrm dA$ exakt
erfüllt sein -- der zentrale Verifikationstest des thermischen Lösers
(siehe Abschnitt~\ref{sec:validation}).


\section{Numerische Umsetzung und Softwarearchitektur}
\label{sec:numerics}

Beide Teilprobleme werden mit der Finite-Elemente-Methode (FEM) in
Galerkin-Formulierung diskretisiert, mit linearen Dreieckselementen (P1)
auf einem strukturierten 2D-Netz $(r,z)$; die achsensymmetrische
Schwachform trägt das Volumengewicht $2\pi r$ (Details in
Anhang~\ref{app:derivation}). Die elektromagnetische Assemblierung ergibt
eine komplexe dünnbesetzte Matrix (wegen $j\omega\sigma$), gelöst mit
\url{scipy.sparse.linalg.spsolve} auf Basis von
NumPy~\cite{harris2020numpy} und SciPy~\cite{virtanen2020scipy}. Bei
$\sim\!10^4$ Unbekannten
(2D, achsensymmetrisch) ist ein direkter Löser ausreichend; die
Skintiefe in Aluminium bei $50\,\mathrm{Hz}$
($\delta\approx12\,\mathrm{mm}\gg3\,\mathrm{mm}$ Scheibendicke) macht eine
feine Vernetzung in $z$-Richtung unnötig.

Aus der einmaligen FEM-Lösung wird ein niederdimensionales,
echtzeitfähiges Modell abgeleitet: eine $I^2$-Skalierung mit
Zeitkonstante erster Ordnung ($\tau\approx4{,}07\,\mathrm{min}$ bei
$R=80\,\mathrm{mm}$) für die Scheibe, sowie ein verteiltes RC-Netzwerk
(\emph{lumped-parameter thermal network}) für Spulen, Eisenkern/-ring und
den gemeinsamen Luftknoten -- ein in der elektrischen Antriebstechnik
etabliertes Modellierungsprinzip~\cite{wallscheid2021thermal}, das sich
direkt an realen Sensordaten kalibrieren lässt. Der gesamte Online-Schritt
reduziert sich auf $O(n)$-Arithmetik, lauffähig in Millisekunden im
Browser intelligenter Endgeräte wie Laptop und Smartphone.
Tabelle~\ref{tab:modules} (Anhang~\ref{app:software})
listet zusätzlich für jede Datei explizit, welche Gleichung sie löst und
mit welcher Bibliothek/Methode.

\begin{figure*}[t]
\centering
\resizebox{\textwidth}{!}{%
\begin{tikzpicture}[
  node distance=5mm and 9mm,
  box/.style={draw, rounded corners, align=center, minimum width=3.0cm,
              minimum height=0.95cm, font=\footnotesize},
  arr/.style={-Stealth, thick}
]
\node[box] (params) {\texttt{params.yaml}\\[-2pt]\tiny (SSOT aller Parameter)};
\node[box, right=of params] (config) {\texttt{config.py}\\[-2pt]\tiny Einheiten, $I_{\mathrm{peak}}$};
\node[box, right=of config] (em) {\texttt{em\_solver.py}\\[-2pt]\tiny EM-FEM, Verlustkarte};
\node[box, right=of em] (th) {\texttt{thermal\_solver.py}\\[-2pt]\tiny Wärme-FEM, Bilanz};
\node[box, right=of th] (rom) {\texttt{rom.py} / \texttt{twin\_core.py}\\[-2pt]\tiny Echtzeit-ROM (SSOT)};
\node[box, right=of rom] (out) {\texttt{build\_twin\_html\_fem.py}\\[-2pt]\tiny Bake $\to$ AR-HTML + QR};
\node[box, below=of rom] (sensor) {\texttt{data\_io.py} + Arduino\\[-2pt]\tiny gemessener Strom $I(t)$};

\draw[arr] (params) -- (config);
\draw[arr] (config) -- (em);
\draw[arr] (em) -- (th);
\draw[arr] (th) -- (rom);
\draw[arr] (rom) -- (out);
\draw[arr] (sensor) -- (rom);
\end{tikzpicture}%
}
\caption{Softwarepipeline: von der Parameterdatei über die
Offline-FEM-Löser zum echtzeitfähigen ROM und der finalen
AR-Auslieferung. Der Sensorpfad liefert den gemessenen Spulenstrom $I(t)$
als \emph{Eingang} des Zeitintegrators (Abschnitt~\ref{sec:realtime_input});
die Temperatur ist die daraus berechnete Ausgangsgröße.}
\label{fig:pipeline}
\end{figure*}

Abbildung~\ref{fig:pipeline} zeigt die Ein-Weg-Pipeline: Jede Datei trägt
genau eine Verantwortung, und alle physikalischen Parameter liegen
zentral in \texttt{params.yaml} (\emph{single source of truth}) -- keine
Konstante ist im Code fest verdrahtet.

\subsection{Gemessener Wechselstrom als Echtzeit-Eingang}
\label{sec:realtime_input}

Bislang wird der Zwilling mit einer \emph{konstanten} Annahme
$I=5\,\mathrm A$ betrieben. Der eigentliche Zielzustand ist ein anderer:
Der Spulenstrom soll \textbf{fortlaufend gemessen} und als Eingangsgröße
in das thermische Modell zurückgeführt werden -- erst damit wird aus der
Simulation ein Zwilling, der dem realen Betrieb folgt statt ihn nur
nachzubilden.

Messtechnisch ist dies gelöst über einen Hall-Effekt-Stromsensor
(ACS712-20A) in Reihe mit einem Zuleiter des Spulenkreises. Er ist
galvanisch getrennt ($\approx2{,}1$\,kV), was bei den anliegenden
$\approx195$--$240\,\mathrm V$ Netzspannung zwingend ist; ein einfacher
Shunt-Widerstand wäre hier nicht zulässig. Der Mikrocontroller bildet den
Effektivwert der Wechselkomponente um den \emph{gemessenen} Mittelwert
(nicht um den nominellen $V_{\mathrm{cc}}/2$-Offset, da sowohl die
$5\,\mathrm V$-Versorgung als auch der Sensor-Nullpunkt Toleranz und Drift
aufweisen):
\begin{equation}
  I_{\mathrm{rms}} = \frac{k}{S}\sqrt{\frac{1}{N}\sum_{n=1}^{N}
      \bigl(v_n - \bar v\bigr)^2},
  \qquad \bar v = \frac{1}{N}\sum_{n=1}^{N} v_n .
  \label{eq:irms}
\end{equation}
Dabei ist $S=0{,}100\,\mathrm{V/A}$ die Sensorempfindlichkeit und $k$ ein
Ein-Punkt-Korrekturfaktor gegen ein Multimeter am
$5\,\mathrm A$-Betriebspunkt ($k=1$, solange nicht gemessen); er absorbiert
die Toleranz beider Datenblattkonstanten.

Entscheidend für die Genauigkeit ist die Wahl des Mittelungsfensters: Ein
echter Effektivwert ist nur über eine \emph{ganzzahlige} Anzahl
Netzperioden exakt. Das Fenster ist deshalb über die \textbf{Zeit}
festgelegt -- fünf Netzperioden, also exakt $100\,\mathrm{ms}$ -- und
nicht über eine feste Stichprobenzahl $N$. Andernfalls hinge das Ergebnis
von der ADC-Geschwindigkeit des jeweiligen Boards ab: dieselben $1000$
Abtastwerte überspannen auf einem AVR-Board etwa $5{,}6$ Netzperioden, auf
dem schnelleren ADC eines Renesas-Boards nur rund $1{,}25$, was einen
phasenabhängigen Fehler von mehreren Prozent verursacht. Mit
zeitbegrenztem Fenster ist $N$ eine Folge der Boardgeschwindigkeit statt
einer Fehlerquelle, und dieselbe Firmware liefert auf beiden
Plattformen denselben Messwert.

Softwareseitig ist die Kette vollständig implementiert: Der Messwert wird
im Sekundentakt als CSV übertragen und treibt den Zeitintegrator direkt an
-- der Zeitschritt wird mit dem gemessenen $I(t)$ statt mit der Konstante
ausgeführt (\texttt{data\_io.py} ruft
\texttt{TwinState.step(}$I_{\mathrm{rms}}$\texttt{,}~$\Delta t$\texttt{)}
mit dem $\Delta t$ zwischen zwei Proben auf; bei ungültiger Probe fällt es
auf den Konstantwert zurück). Da alle Verluste nach
Gl.~\eqref{eq:realtime} quadratisch mit dem Strom skalieren, bleibt es
dabei bei einer skalaren Multiplikation pro Zeitschritt -- die
Echtzeitfähigkeit bleibt also auch mit echtem Messeingang erhalten.
Getestet wurde die Kette durchgängig mit synthetischen Messreihen; es
fehlt allein der physische Einbau des Sensors -- diese Umsetzung ist
derzeit in Arbeit (Stand: 02.09.2026).


\section{Warum Simulation mit Python-Bibliotheken}
\label{sec:tools}

Die Simulation wird bewusst selbst in Python
(NumPy~\cite{harris2020numpy}, SciPy~\cite{virtanen2020scipy})
umgesetzt, weil die eigentliche Anforderung -- ein echtzeitfähiges,
sensor-kalibrierbares, auf einem Smartphone lauffähiges Modell --
Kontrolle auf Quellcode-Ebene voraussetzt: Um die $I^2$-Struktur aus
Abschnitt~\ref{sec:concept} auszunutzen (FEM einmal lösen, danach nur
noch skalieren), muss die Verlustkarte $\hat q(r,z)$ \emph{vor} der
Multiplikation mit $I^2$ gezielt herausgezogen und in einem eigenen
Echtzeit-ROM weiterverwendet werden können -- strukturell nur möglich,
wenn der Löser selbst geschrieben ist.

Zum Einordnen wurden dennoch die etablierten Alternativen geprüft:
Werkzeuge wie SimScale, COMSOL/ANSYS Maxwell~\cite{comsol_team28} oder
FEMM~\cite{meeker2020femm} könnten das elektromagnetisch-thermische
Teilproblem grundsätzlich lösen -- COMSOL bietet für TEAM~28 sogar ein
fertiges Beispielmodell~\cite{comsol_team28} --, ebenso die
Open-Source-Löser Elmer~\cite{malinen2013elmer} und
FEniCS~\cite{langtangen2017fenics}. Tabelle~\ref{tab:tools} zeigt: Keines
davon erfüllt ohne größeren Zusatzaufwand die obige
Echtzeit-Anforderung.

\begin{table*}[t]
\centering\small
\caption{Vergleich der Werkzeugoptionen gegen die Anforderungen des digitalen Zwillings (nicht gegen eine reine Einzelsimulation).}
\label{tab:tools}
% tabularx: X columns split \textwidth exactly and wrap their own text, so the
% table can never overflow the page margin regardless of cell content length
% (plain tabular's c-columns size to content and silently spill past the margin).
\begin{tabularx}{\textwidth}{@{}l
    >{\centering\arraybackslash}X
    >{\centering\arraybackslash}X
    >{\centering\arraybackslash}X
    >{\centering\arraybackslash}X
    >{\centering\arraybackslash}X@{}}
\toprule
Kriterium & SimScale & COMSOL/Maxwell & FEMM & Elmer/FEniCS & \textbf{Python (gewählt)} \\
\midrule
Niederfrequenz-EM (Wirbelstrom) & begrenzt & \checkmark & \checkmark & \checkmark & \checkmark (Verlustpfad ja, Hubkraft offen$^\ast$) \\
EM$\to$Thermik-Kopplung & schwierig & \checkmark & schwach & komplex & \checkmark direkt in-memory \\
ROM / Echtzeitfähigkeit & -- & teilweise, teuer & -- & -- & \checkmark \textbf{Kernentwurf} \\
Export nach Web/AR & -- & -- & -- & -- & \checkmark 1 HTML-Datei \\
Sensor-Kalibrierung & -- & proprietär & -- & selbst zu bauen & \checkmark \url{data_io.py} \\
Läuft auf macOS & \checkmark & \checkmark (teuer) & -- nur Windows & \checkmark & \checkmark \\
Kosten & Core-Hour-Limit & tausende~€ & kostenlos & kostenlos & kostenlos \\
Black Box? & ja & ja & teilweise & nein & \textbf{nein -- volle Kontrolle} \\
\bottomrule
\end{tabularx}

\vspace{2pt}
{\footnotesize $^\ast$Verlustpfad (Wirbelstrom-/Ohmverluste) ist validiert; die
Hubkraft/Levitationshöhe stimmt noch nicht mit der Beobachtung überein
-- siehe Abschnitt~\ref{sec:limitations}.}
\end{table*}

Die letzte Zeile von Tabelle~\ref{tab:tools} fasst den strukturellen
Kern zusammen: Kommerzielle und viele Open-Source-FEM-Pakete kapseln ihren
Löser als Black Box, unabhängig von den übrigen Kriterien -- jede
Stromänderung erfordert dort einen vollständigen Solver-Neustart, statt
der oben beschriebenen einmaligen Lösung plus Skalierung. FEMM entfällt
zusätzlich aus praktischem Grund (nur Windows-nativ, Entwicklungsrechner
ist macOS).

Für einen unabhängigen Vergleichslauf in einem der genannten
kommerziellen Pakete sind Abbildung~\ref{fig:geometry} und
Tabelle~\ref{tab:geometry} (Abschnitt~\ref{sec:intro}) direkt
verwendbar: Sie enthalten dieselben Radien, Windungszahlen und
Materialkonstanten, die auch der eigene Löser aus \texttt{params.yaml}
liest -- ein Nachbau ist damit ohne Rückfrage beim Team möglich.


\section{Validierung und bisherige Ergebnisse}
\label{sec:validation}

\begin{table}[t]
\centering\small
\caption{Verlustleistungen der Offline-FEM bei einer Anregungs\emph{amplitude}
von $5\,\mathrm A$ und $T_\infty=20\,^\circ$C. Zur Umrechnung auf den realen
Betriebspunkt siehe Anmerkung~\ref{rem:current}.}
\label{tab:losses}
\begin{tabular}{@{}lr@{}}
\toprule
Verlustquelle & $P$\,[W] \\
\midrule
Wirbelstromverlust Scheibe (Al) & 25{,}81 \\
Wirbelstromverlust Eisenkern + -ring & 5{,}86 \\
Ohmscher Spulenverlust (innen 52{,}32 + außen 54{,}12) & 106{,}44 \\
\midrule
\textbf{Gesamt} & \textbf{138{,}11} \\
\bottomrule
\end{tabular}
\end{table}

Tabelle~\ref{tab:losses} zeigt die aktuelle Verlustverteilung. Die
folgenden internen Prüfungen sind als Python-Funktionen auf Basis von
NumPy~\cite{harris2020numpy} und SciPy~\cite{virtanen2020scipy}
implementiert und werden bei jeder Modelländerung erneut ausgeführt;
ihre Aussagekraft ist bewusst unterschiedlich gewichtet:

\begin{itemize}[nosep]
  \item \textbf{Energiebilanz: $0{,}000\,\%$ Abweichung.} Am stationären
        Punkt gilt $\int_V q\,\mathrm dV = \oint_\Gamma h(T-T_\infty)\,\mathrm dA$.
        Dies ist die \emph{aussagekräftigste} Prüfung, da sie unabhängig
        assemblierte Volumen- und Randterme gegeneinander stellt.
  \item \textbf{Gebietsgröße: alle Abweichungen $<1\,\%$} (größte:
        $P_{\mathrm{Scheibe}}$ mit $0{,}767\,\%$) beim Vergleich von
        Dirichlet- gegen Neumann-Randbedingung auf dem
        $1\times1\,\mathrm m$-Gebiet -- der vom Fachgebiet vorgegebene
        Nachweis, dass der Außenrand weit genug entfernt liegt.
  \item \textbf{$I^2$-Skalierung: $P(2I)/P(I) = 3{,}994$} (Sollwert
        $4{,}000$) bei vollständiger Neulösung der FEM. Die Restabweichung
        von $0{,}15\,\%$ ist \emph{kein} Fehler, sondern der erwartete
        Effekt der nichtlinearen $\mu_r(B)$-Korrektur: Bei doppeltem Strom
        sinkt $\mu_r$ im Eisen von $993$ auf $977$, sodass die Verluste
        geringfügig unterquadratisch wachsen.
\end{itemize}

\begin{remark}[$I^2$-Prüfung im ROM]
Das ROM selbst liefert für dieselbe Prüfung exakt $4{,}000000$. Dieser
Wert ist jedoch \emph{tautologisch}: Das ROM \emph{ist} per Konstruktion
eine skalare Multiplikation mit $(I/I_{\mathrm{ref}})^2$, kann diesen Test
also gar nicht verfehlen. Er belegt lediglich die korrekte Arithmetik der
Reduktion, nicht die Physik. Aussagekräftig ist allein der oben genannte
FEM-Wert $3{,}994$.
\end{remark}

\begin{remark}[Stromkonvention]
\label{rem:current}
Der am Prüfstand gemessene Wert $I_{\mathrm{rms}}=5\,\mathrm A$ ist ein
Effektivwert; die zugehörige Amplitude beträgt
$\hat I=\sqrt2\,I_{\mathrm{rms}}=7{,}07\,\mathrm A$. Die Verlustkette
(Tabelle~\ref{tab:losses}) setzt bewusst $5\,\mathrm A$ als
Phasor\-amplitude ein. Da alle Verluste $\propto I^2$ sind, liegen die
dort angegebenen Wattzahlen damit um den Faktor~$2$ unter der am realen
Betriebspunkt tatsächlich umgesetzten Leistung. Das ist zulässig und
beabsichtigt, weil dieser konstante Faktor beim Fit der
Konvektionsleitwerte $hA$ gegen die gemessenen Spulentemperaturen
vollständig absorbiert wird -- die \emph{Temperatur}vorhersage bleibt
dadurch korrekt. Die \emph{Kraft}kette rechnet dagegen mit der echten
Amplitude $\hat I$, da eine Kraft nicht wegkalibriert werden kann.
\end{remark}

Die Sättigungsreserve des Eisens ist ebenfalls unkritisch, muss aber
korrekt verglichen werden: Der Löser meldet $B_{\max}=0{,}663\,\mathrm T$
als \emph{Effektivwert}, während $B_{\mathrm{sat}}=1{,}5\,\mathrm T$ eine
Spitzenwert-Materialgröße ist. Der maßgebliche Spitzenwert beträgt somit
$\sqrt2\cdot0{,}663 \approx 0{,}94\,\mathrm T$, also rund $63\,\%$ von
$B_{\mathrm{sat}}$. Die Annahme eines linearen $\mu_r$ bleibt damit
gültig, die Reserve ist jedoch kleiner, als ein direkter Vergleich der
Effektivwerte suggerieren würde.

Zwei reale Betriebspunkte wurden per IR-Kamera vermessen (Variac-Dial
$220^\circ\approx195{,}5\,\mathrm V\to5\,\mathrm A$ als
Hauptbetriebspunkt, Variac-Dial auf Maximalstellung $270^\circ$
($\approx240\,\mathrm V$) $\to7{,}8\,\mathrm A$ als Kalibrieranker). Das
RC-Netzwerkmodell der Spulen wurde direkt gegen den vollständigen,
nichtlinearen Zeitintegrator kalibriert und trifft den stationären
Zustand exakt ($T_{\mathrm{innen}}=79{,}00\,^\circ$C,
$T_{\mathrm{außen}}=74{,}00\,^\circ$C bei $7{,}8\,\mathrm A$). Nur die
dunkel lackierten Spulen liefern verlässliche IR-Werte (unbekannte
Emissivität von blankem Aluminium). Für Scheibe und Kern ist die
naheliegende Abhilfe kein zusätzlicher Sensor, sondern ein
mattschwarzer Messfleck bekannter Emissivität ($\varepsilon\approx0{,}95$)
auf der Scheibenunterseite -- kabellos, und damit ohne Rückwirkung auf die
Levitation.

\textbf{Stationär validiert, transient noch nicht.} Diese Unterscheidung
ist wichtig: Der stationäre Zustand ist per Konstruktion exakt, weil
$hA_{\mathrm{innen/außen}}$ genau darauf gefittet wurden. Der zeitliche
Verlauf ist damit jedoch \emph{nicht} mitvalidiert. Gegen die einzige
vorhandene Transientenmessung -- eine zweistufige Stromrampe
($5\,\mathrm A$ bis $t=300\,\mathrm s$, dann $7{,}75\,\mathrm A$) mit fünf
Zeitstempeln der Außenspulentemperatur -- ergibt das aktuelle Modell einen
Fehler von $\mathrm{RMS}=11{,}5\,^\circ$C (Tabelle~\ref{tab:transient}):
Es erwärmt sich systematisch zu langsam. Ursache ist eine bewusst in Kauf
genommene Modelländerung; sie wird in Abschnitt~\ref{sec:limitations}
diskutiert.

\begin{table}[t]
\centering\small
\caption{Transientenvergleich gegen die gemessene Stromrampe
(Außenspule, $T_\infty=29\,^\circ$C). Der stationäre Zustand stimmt exakt,
der Aufheizverlauf noch nicht.}
\label{tab:transient}
\begin{tabular}{@{}rrrr@{}}
\toprule
$t$\,[s] & gemessen\,[$^\circ$C] & Modell\,[$^\circ$C] & $\Delta$\,[K] \\
\midrule
140 & 40{,}00 & 33{,}55 & $-6{,}45$ \\
200 & 43{,}00 & 34{,}70 & $-8{,}30$ \\
300 & 48{,}50 & 36{,}22 & $-12{,}28$ \\
360 & 55{,}75 & 40{,}00 & $-15{,}75$ \\
450 & 56{,}00 & 43{,}55 & $-12{,}45$ \\
\midrule
\multicolumn{3}{@{}l}{\textbf{RMS}} & \textbf{11{,}53} \\
\bottomrule
\end{tabular}
\end{table}

Zur Rollenverteilung von Messung und Modell sei ausdrücklich festgehalten:
Der Prüfstand besitzt \emph{keine} fest installierte Temperatursensorik.
Gemessen wird ausschließlich der Spulenstrom; die Temperatur ist die
Größe, die das Modell daraus \emph{berechnet}
(Abschnitt~\ref{sec:realtime_input}). Temperaturmesswerte zur Gegenprobe
stammen daher aus manuellen IR-Aufnahmen, nicht aus einem kontinuierlichen
Log. Eine automatisierte Nachkalibrierung des Wärmeübergangs aus einem
Sensorlog ist auf diesem Aufbau folglich nicht möglich -- eine bewusst
akzeptierte Grenze, keine offene Implementierungslücke.


\section{Erwartete Ergebnisse und Ausblick}
\label{sec:outlook}

Kurzfristig erwartet: die Inbetriebnahme der Sensorhardware und damit der
Übergang von der konstanten Annahme $I=5\,\mathrm A$ auf einen
\textbf{real gemessenen Stromverlauf $I(t)$}, der den thermischen Zwilling
zur Laufzeit speist (Abschnitt~\ref{sec:realtime_input}); daraus eine reale
Aufheiz-\emph{und} Abkühlkurve, die das bislang unteridentifizierte
RC-Netzwerk (insbesondere den gemeinsamen Luftknoten und die
Zwei-Knoten-Aufteilung der Spule) erstmals identifizierbar macht und damit
die derzeit offene Transientenabweichung (Abschnitt~\ref{sec:limitations})
schließen soll -- die Temperaturseite dieser Kurven wird dabei
berührungslos per IR-Kamera erfasst, mit mattschwarzem Messfleck auf der
Scheibenunterseite als Abhilfe gegen die Emissivitätsunsicherheit; sowie
eine veröffentlichte AR-Ansicht mit QR-Code, aufrufbar ohne Installation. Mittelfristig: eine dichtere Kalibriertabelle
Stelldrehwinkel~$\to$~Strom und eine belastbarere Bestimmung der
Eisenpermeabilität zur Klärung der in Abschnitt~\ref{sec:limitations}
genannten Hubkraft-Frage.


\section{Offene Fragen und Grenzen}
\label{sec:limitations}

Im Sinne wissenschaftlicher Transparenz: Die vorhergesagte
Levitationshöhe ($z_{\mathrm{eq}}\approx11{,}7$--$14{,}7\,\mathrm{mm}$)
stimmt nicht mit dem beobachteten Spalt ($7$--$8\,\mathrm{mm}$) überein;
magnetische Sättigung wurde als Ursache ausgeschlossen (Kraftänderung
$<0{,}1\,\%$), wahrscheinlichste Ursache ist der unvermessene
$\mu_r$-Platzhalterwert des Eisens. Eine willkürliche Rückwärtsanpassung
von $\mu_r$ an die beobachtete Höhe wurde bewusst \emph{nicht}
vorgenommen. Der Original-TEAM-28-Benchmark (ohne Eisen, $20\,\mathrm A$)
weicht ebenfalls $\approx28\,\%$ von der Referenztabelle ab -- pausiert,
nicht blockierend.

\textbf{Transientenabweichung der Spule.} Der in
Tabelle~\ref{tab:transient} dokumentierte Fehler von
$\mathrm{RMS}=11{,}5\,^\circ$C ist die derzeit größte offene Diskrepanz
des thermischen Modells. Ihre Ursache ist bekannt und nachvollziehbar: Die
Spule wird als Zwei-Knoten-Körper modelliert (IR-sichtbare Oberfläche plus
träges Wicklungsinneres). In einer früheren Modellversion wurde die
gesamte ohmsche Verlustleistung auf den Oberflächenknoten gelegt, der nur
$22{,}4\,\%$ der Wärmekapazität trägt -- das ergab eine gut zur Messung
passende, aber physikalisch unbegründete Aufheizrate. Die Verlustleistung
wird inzwischen massenproportional auf beide Knoten verteilt, was den
stationären Zustand beweisbar unverändert lässt, die anfängliche
Aufheizrate jedoch um den Faktor $\approx4{,}5$ senkt. Das Modell ist
seither physikalisch konsistenter, trifft die gemessene Rampe aber
schlechter. Die dafür maßgeblichen Parameter (Kapazitätsaufteilung
\texttt{coil\_C\_scale}, innere Kopplung \texttt{coil\_G\_wind}) sind aus
einer reinen Aufheizmessung nicht eindeutig bestimmbar -- sie erfordern
zusätzlich eine \emph{Abkühlkurve}, die noch aussteht. Bis dahin gilt: der
stationäre Zustand ist validiert, der Zeitverlauf ausdrücklich nicht.

Ohne diese Abkühlkurve bleibt zudem der gemeinsame
Luftknoten des RC-Netzwerks unteridentifiziert. Ebenfalls offen: Im
Modell erreichen Scheibe und Spule annähernd dieselbe Temperatur, während
die IR-Daten eine deutlich heißere Spule nahelegen; entschieden werden
kann das erst mit einer vertrauenswürdigen Ablesung an der
Scheibenunterseite, also mit dem oben genannten mattschwarzen Messfleck.
Schließlich ist die $\sigma(T)$-Rückkopplung der Spule
derzeit nur implizit über den $hA$-Fit enthalten
(Anmerkung~\ref{rem:sigma}).


\section{Zusammenfassung}
\label{sec:conclusion}

Das Projekt verfolgt einen physikalisch begründeten Weg zu einem
echtzeitfähigen digitalen Zwilling: Eine einmalige, hochauflösende
FEM-Lösung wird über $I^2$-Skalierung und ein kalibrierbares
Reduced-Order-Model auf Millisekunden-Laufzeit reduziert -- eine Struktur,
die mit gekapselten kommerziellen CAE-Werkzeugen nur schwer umsetzbar
ist. Jede Änderung am Modell wird ausschließlich über harte,
automatisierte physikalische Kriterien freigegeben (Energiebilanz,
$I^2$-Konsistenz). Erste
Kalibrierungsergebnisse an der Spulentemperatur bestätigen die
Tragfähigkeit des Modells; die Levitationshöhe bleibt als offene
physikalische Frage bestehen und wird mit der geplanten Sensorhardware
weiter untersucht.


%%%%%%%%%%%%%%%%%%%%%%%%%%%%%%%%%%%%%%%%%%%%%%%%%%%%%%%%%%%%%%%%%%%%%%%%%%%%%%%%
\appendices

\section{Vollständige Herleitung der physikalischen Gleichungen}
\label{app:derivation}

Die axisymmetrische Schwachform der Wärmeleitungsgleichung mit
Volumengewicht $2\pi r$ lautet
\begin{multline*}
  \int k\,\nabla T\!\cdot\!\nabla v \cdot 2\pi r \,\mathrm dA
  + \oint_\Gamma h\,T\,v\cdot 2\pi r\,\mathrm ds \\
  = \int p\,v\cdot 2\pi r\,\mathrm dA
  + \oint_\Gamma h\,T_\infty\,v\cdot 2\pi r\,\mathrm ds .
\end{multline*}
Für P1-Elemente ist $\nabla T$ elementweise konstant, sodass die
Elementsteifigkeitsmatrix
$K_e = 2\pi k\,(bb^\top+cc^\top)\,A\,r_c$ direkt aus den Formfunktionsgradienten
$b,c$, der Elementfläche $A$ und dem Schwerpunktradius $r_c$ folgt. Die
Konvektionsrandmatrix mit linearer Radiusgewichtung ergibt sich zu
$K_{\mathrm{edge}} = 2\pi h\,\tfrac{L}{12}
\left(\begin{smallmatrix} 3r_a+r_b & r_a+r_b \\ r_a+r_b & r_a+3r_b\end{smallmatrix}\right)$.
Im \emph{thermischen} Teilproblem ist die Achse $r=0$ eine
Symmetriebedingung und bleibt als natürliche Randbedingung unerzwungen
(Randkanten auf der Achse werden bei der Konvektionsassemblierung
übersprungen). Im \emph{elektromagnetischen} Teilproblem gilt das
\textbf{nicht}: Dort wird auf der Achse $A_\varphi=0$ als
Dirichlet-Bedingung explizit gesetzt, da die azimutale Komponente des
Vektorpotentials bei $r\to0$ physikalisch verschwinden muss. Für dieses
Teilproblem folgt analog eine
Steifigkeitsmatrix aus Gl.~\eqref{eq:em}. Der Term $j\omega\sigma$ macht
sie komplexwertig; sie wird als \emph{eine} komplexe Matrix assembliert
und direkt in komplexer Arithmetik faktorisiert (keine Aufspaltung in
Real- und Imaginärteilsystem). Die magnetische Flussdichte ergibt sich
aus $B_r=-\partial A_\varphi/\partial z$,
$B_z=\tfrac1r\,\partial(rA_\varphi)/\partial r$.

Physikalisch beschreibt Gl.~\eqref{eq:em} das AC-Feld: Der Term $j\omega\sigma
A_\varphi$ ist die Wirbelstromreaktion in jedem leitfähigen Körper
(Scheibe, Eisen); der Term $\nu A_\varphi/r^2$ ist der geometrische Effekt
der Achsensymmetrie und der Grund, warum 2D statt 3D genügt. Bei
$I_{\mathrm{rms}}=5\,\mathrm A$ durch die $1000$-Windungen-Innenspule
beträgt $B_{\max}$ im Eisenkern $0{,}663\,\mathrm T$ als Effektivwert,
entsprechend einem Spitzenwert von $\approx0{,}94\,\mathrm T$ -- gegenüber
$B_{\mathrm{sat}}=1{,}5\,\mathrm T$ also unsaturiert (vgl.
Abschnitt~\ref{sec:validation}).

\section{Vollständige Softwarepipeline}
\label{app:software}

Tabelle~\ref{tab:modules} beantwortet explizit, welche Bibliothek welche
Gleichung mit welcher numerischen Methode löst -- ergänzend zu
Abbildung~\ref{fig:pipeline} (Datenfluss) und Anhang~\ref{app:derivation}
(Herleitung der Gleichungen selbst).

\begin{table}[h]
\centering\scriptsize
\caption{Kernmodule der Softwarepipeline: gelöste Gleichung und
numerische Methode je Datei.}
\label{tab:modules}
\begin{tabularx}{\linewidth}{@{}p{1.6cm} X X@{}}
\toprule
Datei & Löst & Methode \\
\midrule
\url{config.py} & Parameter laden/normieren (mm$\to$m), $I_{\mathrm{peak}}=I_{\mathrm{rms}}\sqrt2$ & PyYAML; Datenklassen, kein Gleichungslöser \\
\url{em_solver.py} & Gl.~\eqref{eq:em} $\to$ Gl.~\eqref{eq:loss}, Hubkraft & \url{scipy.sparse.linalg.spsolve} (komplex, LU) \\
\url{thermal_solver.py} & Gl.~\eqref{eq:fourier}, stationär & \url{scipy.sparse.linalg.spsolve} (reell, SPD) \\
\url{rom.py} & $\tau{=}C/UA$ aus $\hat P,\Delta\hat T$; Verifikations-ODE & \url{scipy.integrate.solve_ivp} \\
\url{twin_core.py} & RC-Netzwerk (Spule/Eisen/Luft); Levitation & explizites Euler, $n_{\mathrm{sub}}{=}\lceil\Delta t/(0{,}05\tau)\rceil$; Levitation geschlossen gelöst \\
\url{refit_hA.py} & Nullstelle von $T_{\mathrm{ss}}(hA)-T_{\mathrm{meas}}$ & \url{scipy.optimize.fsolve} \\
\url{xval_twin.py} & -- (Software-Verifikation) & Playwright: Python- vs.\ JS-Integrator, $|\Delta|{<}10^{-9}$ \\
\url{build_twin_html_fem.py} & bake: identische Formeln in JavaScript & -- \\
\url{data_io.py} & Sensorbrücke $\to$ \url{TwinState.step(I,dt)} & seriell/Mock-CSV-Parser (NumPy) \\
\bottomrule
\end{tabularx}
\end{table}

Algorithmisch gliedert sich die Pipeline exakt in die zwei Stufen aus
Abschnitt~\ref{sec:concept}: \textbf{Offline} lösen \url{em_solver.py}
und \url{thermal_solver.py} je ein lineares Gleichungssystem
\emph{einmal direkt} (LU-Zerlegung, keine Iteration) und liefern
$\hat q(r,z)$, $\hat P_{\mathrm{Cu}}$ sowie $\tau$. \textbf{Online}
integriert \url{twin_core.py} das resultierende RC-Netzwerk mit einem
handgeschriebenen expliziten Euler-Verfahren -- gewählt, weil erst die
Bibliotheksfreiheit (\texttt{numpy+stdlib only}) den späteren
Browser-Export nach JavaScript ermöglicht; die Schrittweite wird pro
Aufruf automatisch in $n_{\mathrm{sub}}$ Teilschritte von höchstens
$5\,\%$ von $\tau$ geteilt, damit das Verfahren auch bei großen
$\Delta t$ (z.\,B. Zeitraffer) numerisch stabil bleibt. Einzige Ausnahme
ist die Levitationsdynamik: Ihre Bewegungsgleichung (gedämpfter
Feder-Masse-Schwinger) besitzt eine exakte analytische Lösung und wird
deshalb bei jedem Aufruf direkt ausgewertet, nicht iterativ integriert.

\bibliographystyle{IEEEtran}
\bibliography{dt4tm_references}

\end{document}
